\section{Описание портов}

Со стороны процессора кэш --- ведомое устройство AXI4-Lite, со стороны внешней памяти --- ведущее.

\begin{table}[H]
	\centering
	\label{tab:cache_ports}
	\small
	\begin{tabularx}{\textwidth}{l c c X}
		\toprule
		\textbf{Порты} & \textbf{Направление} & \textbf{Разрядность} & \textbf{Описание} \\
		\midrule
		
		\verb|clk_i| &
		Вход &
		\verb|1| &
		Тактовый сигнал кэша. \\
		
		\verb|rstn_i| &
		Вход &
		\verb|1| &
		Сигнал синхронного сброса с активным низким уровнем. \\
		
		\midrule
		\multicolumn{4}{l}{\textbf{Интерфейс AXI4-Lite со стороны CPU}} \\
		\midrule
		
		\verb|s_axi_awaddr_i| &
		Вход &
		\verb|ADDR_WIDTH| &
		Адрес записи. \\
		
		\verb|s_axi_awvalid_i| &
		Вход &
		\verb|1| &
		Достоверность адреса записи. \\
		
		\verb|s_axi_awready_o| &
		Выход &
		\verb|1| &
		Готовность принять адрес записи. \\
		
		\verb|s_axi_wdata_i| &
		Вход &
		\verb|DATA_WIDTH| &
		Данные записи. \\
		
		\verb|s_axi_wstrb_i| &
		Вход &
		\verb|DATA_WIDTH / 8| &
		Маска достоверных байтов записи. \\
		
		\verb|s_axi_wvalid_i| &
		Вход &
		\verb|1| &
		Достоверность данных записи. \\
		
		\verb|s_axi_wready_o| &
		Выход &
		\verb|1| &
		Готовность принять данные записи. \\
		
		\verb|s_axi_bresp_o| &
		Выход &
		\verb|2| &
		Ответ на запись. \\
		
		\verb|s_axi_bvalid_o| &
		Выход &
		\verb|1| &
		Достоверность ответа на запись. \\
		
		\verb|s_axi_bready_i| &
		Вход &
		\verb|1| &
		Готовность принять ответ на запись. \\
		
		\verb|s_axi_araddr_i| &
		Вход &
		\verb|ADDR_WIDTH| &
		Адрес чтения. \\
		
		\verb|s_axi_arvalid_i| &
		Вход &
		\verb|1| &
		Достоверность адреса чтения. \\
		
		\verb|s_axi_arready_o| &
		Выход &
		\verb|1| &
		Готовность принять адрес чтения. \\
		
		\verb|s_axi_rdata_o| &
		Выход &
		\verb|DATA_WIDTH| &
		Данные чтения. \\
		
		\verb|s_axi_rresp_o| &
		Выход &
		\verb|2| &
		Ответ на чтение. \\
		
		\verb|s_axi_rvalid_o| &
		Выход &
		\verb|1| &
		Достоверность данных чтения. \\
		
		\verb|s_axi_rready_i| &
		Вход &
		\verb|1| &
		Готовность принять данные чтения. \\
		
		\midrule
		\multicolumn{4}{l}{\textbf{Интерфейс AXI4-Lite со стороны внешней памяти}} \\
		\midrule
		
		\verb|m_axi_awaddr_o| &
		Выход &
		\verb|ADDR_WIDTH| &
		Адрес записи. \\
		
		\verb|m_axi_awvalid_o| &
		Выход &
		\verb|1| &
		Достоверность адреса записи. \\
		
		\verb|m_axi_awready_i| &
		Вход &
		\verb|1| &
		Готовность принять адрес записи. \\
		
		\verb|m_axi_wdata_o| &
		Выход &
		\verb|DATA_WIDTH| &
		Данные записи. \\
		
		\verb|m_axi_wstrb_o| &
		Выход &
		\verb|DATA_WIDTH / 8| &
		Маска достоверных байтов записи. \\
		
		\verb|m_axi_wvalid_o| &
		Выход &
		\verb|1| &
		Достоверность данных записи. \\
		
		\verb|m_axi_wready_i| &
		Вход &
		\verb|1| &
		Готовность принять данные записи. \\
		
		\verb|m_axi_bresp_i| &
		Вход &
		\verb|2| &
		Ответ на запись. \\
		
		\verb|m_axi_bvalid_i| &
		Вход &
		\verb|1| &
		Достоверность ответа на запись. \\
		
		\verb|m_axi_bready_o| &
		Выход &
		\verb|1| &
		Готовность принять ответ на запись. \\
		
		\verb|m_axi_araddr_o| &
		Выход &
		\verb|ADDR_WIDTH| &
		Адрес чтения. \\
		
		\verb|m_axi_arvalid_o| &
		Выход &
		\verb|1| &
		Достоверность адреса чтения. \\
		
		\verb|m_axi_arready_i| &
		Вход &
		\verb|1| &
		Готовность принять адрес чтения. \\
		
		\verb|m_axi_rdata_i| &
		Вход &
		\verb|DATA_WIDTH| &
		Данные чтения. \\
		
		\verb|m_axi_rresp_i| &
		Вход &
		\verb|2| &
		Ответ на чтение. \\
		
		\verb|m_axi_rvalid_i| &
		Вход &
		\verb|1| &
		Достоверность данных чтения. \\
		
		\verb|m_axi_rready_o| &
		Выход &
		\verb|1| &
		Готовность принять данные чтения. \\
		
		\bottomrule
	\end{tabularx}
	\caption{Порты кэша}
\end{table}
